% BEGIN_REPORT::LEAD::ARC-0002::2026-10-07T15:59:43-07:00
\documentclass[11pt]{article}
\usepackage[margin=1in]{geometry}
\usepackage{amsmath,amssymb,amsthm,booktabs,hyperref}
\newtheorem{theorem}{Theorem}
\newcommand{\SAT}{\mathrm{SAT}}
\newcommand{\R}{\mathcal R_{\SAT}}
\newcommand{\B}{\mathcal B}
\title{ARC-0002: Cofinal seed hitting as a weaker sufficient hypothesis}
\author{LEAD}
\date{7 October 2026}
\begin{document}
\maketitle

\paragraph{Result and status.}
\textbf{PROVED/EXACT logical implication; SAT-specific premise OPEN.}
The contradiction argument does not require one improvement law on all of
$\R$, nor a common contraction factor.  It is enough to prove finite hitting
from every achievable member of one componentwise cofinal seed family.  The
diagonal seeds $(K,K)$ for powers of two already form such a family.  This
reduces the quantifiers of the open arrow but supplies no SAT route.  Independent
RED review of the new statement remains pending, and $P$ versus $NP$ remains
open.

\section*{Exact claim card}
Work in the fixed deterministic random-access model and explicit SAT encoding
required by the Williams endpoint.  The achievable region $\R$ contains pairs
whose time and space bounds are witnessed jointly by one uniform exact SAT
machine; it is upward closed.  Let
\[
 \B=\{(t,s)\in\mathbb R_{\ge0}^2:t+s<2\cos(\pi/7)\}.
\]
A set $H\subseteq\mathbb R_{\ge0}^2$ is componentwise cofinal when
$\forall z\,\exists h\in H\;(z\le h)$.

\begin{theorem}[COFINAL-SEED-HITTING1]
If there is a componentwise cofinal $H$ such that
\[
 \forall h\in H\quad
 h\in\R\Longrightarrow
 \exists m<\infty\ \exists z^{(0)},\ldots,z^{(m)}
 \quad z^{(0)}=h,\quad
 (\forall i\le m\ z^{(i)}\in\R),\quad z^{(m)}\in\B,
\]
then $P\ne NP$, subject to exact endpoint-model matching.
\end{theorem}
\begin{proof}
Assume $P=NP$.  Some finite $z$ is in $\R$.  Cofinality gives $h\in H$ with
$z\le h$, and upward closure gives $h\in\R$.  The premise yields a finite
route ending in $\R\cap\B$, contradicting the endpoint theorem
\cite{williams2010}.
\end{proof}

Taking $H=\{(K,K):K\in S\}$ for any unbounded $S\subseteq\mathbb N$ is
sufficient.  Under $P=NP$, the achievable machine and its finite exponent pair
are fixed first; one fixed $K$ is then selected, followed by one fixed finite
route.  Neither the map from $K$ to a route nor the route length must be
uniform over $K$.  In particular, no universal compiler from arbitrary source
code is used.  Conversely, a cap $K(N)$ or route depth growing with input
length is a different, uniform-algorithm claim and is not covered.

\paragraph{Same-realization obligation.}
Every $z^{(i)}\in\R$ asserts the existence of one uniform exact machine with
both coordinates.  A time coordinate from one machine and a space coordinate
from another do not form a route point.  Different route points may have
different witnesses because only finitely many existential consequences are
used after the seed has been fixed.

\section*{Arithmetic model and falsification controls}
For $S=\{2^j:j\ge1\}$, define a seed-dependent arithmetic successor
$z\mapsto\lambda_Kz$ at seed $(K,K)$, where
$\lambda_K=(K-1)/K$.  Every fixed route reaches $t+s<1$ in finitely many
steps, although $\sup_K\lambda_K=1$ and the route lengths are unbounded.  This
demonstrates why a common $\lambda<1$ is unnecessary for the logical
implication; it is not a SAT successor construction.

The exact-rational test
\path{experiments/proof_spines/cofinal_seed_hitting_test.py} ran with seed
20261007, a 25-second cooperative alarm, and elapsed time 0.020691 seconds.
For $K=2,4,8,16,32,64,128,256$, the verified route lengths are respectively
$3,8,21,54,131,309,708,1594$.  Sixty-four seeded rational points were padded
to dominating power-of-two caps; the largest chosen cap was 8192.  A bounded
seed set ending at 16 failed on the point 17.  The current-scale control
$\lambda=1$ failed to cross after 2048 iterations.  The witness-identity check
accepted joint coordinates and rejected coordinate borrowing.  Raw output is
\path{results/raw/proof_spines/2026-10-07_155943_cofinal_seed_hitting.json};
the source digest is
\texttt{ddb6b45f28b972861554ca757ce1dcdbcf4321b3437fab665225104b229d4dd2}.

The test does not infer asymptotics from these values.  Easy SAT families do
not supply the universal route premise.  Padding is used only through upward
closure and cannot erase an unresolved current-scale call.  Independent
products and delayed/source-carry encodings likewise do not imply a route.
The logical implication relativizes to any setting having the two abstract
premises, but the cited forbidden endpoint is not thereby known to relativize
or algebrize.  DEAD, COPY and DELAY transforms can therefore attack a proposed
SAT route even though they do not refute the abstract implication.

\section*{Sources and transfer limits}
The run read the current root \texttt{main.tex}, \texttt{main2.tex},
\texttt{main3.tex}, \texttt{main4.tex}, \texttt{main5.tex}, governance files,
the proof-spines protocol, current peer reports, cited raw experiments, and
the original double-wrapper implementation.  Repository searches found no
prior cofinal-seed theorem.  A connected Drive search for ARC diagonal seed
closure found no corresponding proof artifact; unrelated historical records
were not treated as evidence.

Williams's 2010 theorem supplies the model-sensitive forbidden endpoint
\cite{williams2010}.  Williams's 2024 Circuit-SAT self-improvement gives a
representation-specific finite bootstrapping precedent, but does not prove
ordinary-SAT cofinal routes under $P=NP$ \cite{williams2024}.  Common
polynomial Lyapunov functions for arbitrary switching are an analogy for
uniform contraction requirements, not a transfer to existential SAT witnesses
\cite{mason2004}.  The new theorem is elementary order logic and depends on
none of those analogies beyond the separately cited endpoint.

\section*{Manuscript reconciliation and handoff}
Root \texttt{main4.tex} now contains the theorem, proof, diagonal corollary,
quantifier boundary and arithmetic countermodel to the necessity of a common
contraction.  Root \texttt{main5.tex} receives a dated review stating that no
WRAPPER claim changes.  No OPEN arrow is promoted.

The next target for PROOF and RESEARCH is deliberately narrower: seek a
SAT-specific finite route only from diagonal caps in a sparse unbounded set,
for example powers of two.  RED should try to break such a candidate with
source carry, current-scale calls, padding, products and oracle/algebraic
replacement, and must reject any proof whose cap or depth depends on the input
length.  A single valid cofinal-family route theorem would suffice; reproducing
global ARC would be unnecessary.

\bibliographystyle{plain}
\bibliography{spines}
% END_OF_VERIFIED_REPORT::LEAD::ARC-0002
\end{document}
